\chapter{湍流}
本章讲述不可压缩流体运动的稳定性分析方法和充分发展瑞流的统计理论，介 绍常用的瑞流统计模式及其在工程中的应用.本章最后简要地介绍湍流相干结构的 现，象.

\section*{9.1  湍流的发生}\addcontentsline{toc}{section}{9.1 湍流的发生}

第 8 章中研究了不可压缩牛顿流体运动的基本特性，导出了它的控制方程和边 界条件，并且在若干初边值条件下得到了确定性 解，例如哈根泊扇叶流动等.但是实际流动巾， 只有当无量纲参数 Re 在某一限定值以下，这些 解才能描述真实的流体运动，当 Re 超过某 一 限 定值后.真实流动具有非常复杂的不规则运动状 态.这种运动状态称为湍流或紊流.T.程和自然 界中遇到的流动现象绝大多数是湍流.因此有必 要对它进行研究. 染色些1"- 首先从湍流发生的过程来认识湍流和层流 $\nu$F 的本质差别.雷诺<1 883 )最早用流动显示方法观 k 察团管中流动的湍流发牛.过程.他的实验装置相 --寸工 当简单(见图 9. 门，清水从一个有恒定水位的水 !到 9.1 雷诺实验装置

箱流入等截面直图管.在图管人口的中心处通过一细针孔注入有色液体，以观察团管 内的流动状态.在回管出口处有一阀门调节流量，以改变流动的 R f'.为了减少入口的 扰动，常将入 IJ 磨图成钟罩形. 实验时可用容积法测量流量 Q. 管内的平均流速及 Re 定义为 U... = \_i豆

\begin{gather*}
\pi d ~\\
Rf' = Urnd f'-- -
\end{gather*}

(9. 1) (9.2) 其中 d 为管径， $\nu$ 为水的运动粘性系数.试验过程中，逐渐开大阀门，这时管内流速逐 渐增大.开始时中心染色线保持平稳的直线 状态(图 9. 2( a 门，当流动达到某 - R f'时染 色线开始出现波形扰动(图 9.2(b)). 继续 增大流桂时染色线由剧烈振荡到破裂，并很 快和清水掺挝以致不同能分辨出染色液线 (图 9. 2(c)). 上述第-阶段流动状态称为层 流，最后阶段的流动状态称为湍流. JIJ 间阶段 的流动状态很不稳定，称为过渡流动. 在不加特殊控制的情况下，因管内出现 湍流状态的 Re 约为 2300. 这 一 实验数据常 作为工程管路计算中判断流动状态的依据. (a) (...i'--""'"产\_\_\_\_\_\_\_\_，..飞、j (b) (C) .... .，"""川节，川 斗、二).-...." . 川.... .. ".- ....... ..., ...\textasciitilde{}.. ,,- :' \_. :、.. '!\textasciitilde{}. 飞

\["... -" :. , .. ' ......~'-"\]

、 . .1 ".". 在特殊控制环境下，使外界扰动非常微弱，管 图 9.2 雷诺数增加时流动状态的演变 $\cdot $渡状态; (c) 接近 1;101 流状态 内流动的层流状态可维持到 Re= ](f'. 这就 是说从层流发展到湍流的过程受环境影响很大.因此这 一 阶段流动的性质是极其复 杂的.阿管中流动由民流过渡到湍流的现象在其他流动中(如:边界层、自由剪切后 流动)也能观察到，不过在不同的流动中流动过渡过辑的形态相差很大. 用近代的测速仪器，如热膜流速计、激光测速仪等测量边界层流动中某点流速的 时间序列，则由它的频谱(速度信号时间序列的傅里叶分析) Jt，可以看到，除了在流动 中不可避免的微弱背景扰动外.当 Re 达到某一限定值时出现某种频率的扰动 〈图 9.3(b)) ，继之又产生多种频率成分(图 9.:Hc)) ，随着频率成分的增多.在每一峰 值处频带加宽，直至最终成为具有宽频带连续谱的扰动(图9.3(d)). 图 9. 3 (川， 图9.3( 切，图 9. 3(c). 图9.3(d) 的上部为速度的时间序列.下部为对应的频谱. 从以上直观的观察和时间序列分析中可以了解到从层流到湍流发展的基本过 程为 z (1)层流:它是符合纳维 - 斯托克斯方程在给定初边值条件下的确定性解; (2) 过渡过程的初始阶段 z 出现时间上(或空间上，或同时在空间和时间上)的

l\textasciitilde{}"\textasciitilde{} 川 \{$\omega$1卫生二(ru)Ll\_\_ $\cdot $$\omega$' k- "$\omega$) 队 (a) \} ku (c) (d) !制 9.3 t由流发生过程的 1 J.t问序列和频 i普 (a) )\textasciitilde{}流; (b) 过渡初始阶段; ((")过流发展阶段; Cd) 湍流 周期性扰动 z (:3) 过渡过程的发展:由现多种周期的窄带扰动，现则性流动逐渐破坏 z (4) 湍流 z 含有宽颇带连续谱扰动的完全不规则流动. 从近代非线性动力学系统的角度来看，不可 fk 缩牛顿流体的运动属于含参量 低e) 的非线性艳散系统，层流运动是这 一 系统的确定性解.而湍流是这一系统在确 定条件下的非确定性解.也就是说后流和湍流是这 一 动力学系统的两种状态.究竟发 生哪一种流动状态依赖于动力系统巾的控制参量 : Re.

\section*{9.2  流动的稳定性分析}\addcontentsline{toc}{section}{9.2 流动的稳定性分析}

\section*{9.1  节描述的湍流发牛. 过程是极具复杂的.目前还缺少有效的分析方法来预测它}\addcontentsline{toc}{section}{9.1 节描述的湍流发牛. 过程是极具复杂的.目前还缺少有效的分析方法来预测它}

的全部过程，但是对过踵的初始阶段可以用线性稳定性理论作定量的分析.这 一 分析可 以用来判断什么条件下开始发生层流向湍流过渡，因此在实用上也是很有意义的. 过渡过程的初始阶段特征是，?吁流动参数(如 Re) 增大到一定值后，在民流的确 定性解上将产吃不能自行衰退的周期性扰动.稳定性理论能够确定发生这种扰动的 临界参数以及产生的扰动形态.如频率、波氏、扰动增伏率等.下面介绍不可月- 缩牛顿 流体运动稳定性分析方法及平面泊南叶流动稳定性分析结果. 1.稳定性问题的提法 不可压缩牛顿型?在体运动满足纳维-斯托克斯方程(无量纲形式) : r7 P , 1 r) 2 U , 寻寻 = I， 一了一 十 r 了一了一 Ul <1.1., l"l e O .I，'人 l. ， (9.3a)

aU. L二.L =O

\[<lIj\]

(9.3b) 求解该方程的边界条件和起始条件是 z 在流动边界 2 上

\begin{equation*}U, = F,(x. t)\tag{9.4}\end{equation*}

初始时刻在流场 D 中

\[Ui(X.O) = Gi(x)\]

式中 Fi(x.t) 为在界面 Z 上的给定速度 .G， (x) 为初始速度场. 方程式 (9.3 川和方程式 (9.3b) 在初始值 (9. 5) 和边界条件( 9. 的下的确定性解 U, (x. t) 称为基本流动，一般来说 .U ， (x. t) 可以是定常的，也可以是非定常的，本书只 讨论定常基本流的稳定性问题. 当流动的边值条件不变而初值偏离基本流 Ui(x.t) 的初值 H才.将产生一种新的 (9.5) 流动 U: (x. t). 它满足以下的初边值. 在流动边界 Z 上

\[U~(I.t) = Fi(.r ,t) (9. 6)\]

初始时刻在流场 D 中

\[U~(X.O) = G~(x)\]

将新的流动解口'(x. t) 与基本流 Ui(x. t) 之差用以下公式表式 z

\begin{gather*}
J, (x) = G:(x) - G, (x)\tag{9.8a}\\
U ,(x. t) = U:(x. t) - Ui(x ,t)\tag{9.8b}
\end{gather*}

J\{ i 称为初始扰动 ， Uj 称为流场扰动.稳定性理论将研究扰动 Ui (x.t) 的发展性质. 判断流动稳定性的条件是:如果在任意初始扰动下，基本流的扰动解 Ui (x ，t) 在 时间上是渐近衰减的.那么基本流是稳定的，否则是不稳定的. (9. 7) 2. 扰动方程及线化扰动方程 要由扰动的解来判断基本流的稳定性，必须导出扰动满足的方程及初边值条件. 基本流 Ui 及受扰流动以都满足纳维-斯托克斯方程，即它们分别满足以下方程 z aUi I TT au, a P I 1 a2 U $\tau$\textasciitilde{}+Uj $\tau$一=一了一+古了$\tau$$\iota$ +!i r1 C t1Ij r1.ri 1((;' r1Ij t1Ij au, \{\}Ii 以及 4J + U-h 俨一句

\begin{gather*}
l-b +\\
JW-h
\end{gather*}

一一 川-hUo += 以万川

\[-h\]

将以上两方程相减得扰动方程为 但L 十 U; 丛 +u ， 坐\textasciitilde{} + u ， 主!...=一旦旦+土之\textasciitilde{} (9.9a) ,11 ' -, a.T , ' -, (J.1', ' --, rJ X , (J X;' Re (J Xj (J.Tj

\begin{equation*}eJ u,\tag{9.9b}\end{equation*}

${\rm e}$J x; 其巾 Uj= 口'-U; 和 $\rho$ =p'-p 分别是扰动速度场和扰动压强场，同理将 U; ， 以相应 的初边值方程相减，得 在流动边界 Z 上

\[u, = 0 (9. 10)\]

在流场 D 中

\begin{equation*}u;(x.O) = /;(x)\tag{9.11}\end{equation*}

可以看到扰动发展方程是齐次非线性方程在齐次边界下的初值问题.一般情况 下非线性项 u ，${\rm e}$J ll ，/，1 X， 给定性的分析和数值求解带来很大困难，物理上来说也是由 于这一非线性效应与扩散效应的藕合使过搜过程成为非常复杂的流动.在过渡过程 的起始阶段，由于扰动很小，因而可以用线性化的方法使扰动方程大为简化.假定 U:$\tilde{U}$, \textasciitilde{} 0 (1) (9. 12a)

\begin{equation*}g , \tilde{u}, ~ ()(E)\tag{9.12b}\end{equation*}

其"， :$\epsilon$<< 1.这时扰动方程中的非线性二阶项可以略去，得到线化的扰动方程如下 z -r u-a J-TA --nm 十鱼缸 一一 旦 hu + 吨叫-hUo += h-h 忧-h (9.13a) (9.13b) 下面将利用线化扰动方程 (9.13) 及其初边值公式 (9.10) 和 (9. 11) 讨论稳定性问题. 3. 线性稳定性分析和临界参量 方程(!). 13 川，方程 (9.13b) 是线性的.它们有对时间变量 t 和空间变量的可分离 变:ht解: u; (x ,t) = 玩 ;(x)cxp(- i$\omega$ t)+c.c. (9. 14a)

\begin{equation*}\rho  (x.t) = \theta (x)cxp(- i\omega  I)+c.c.\tag{9.14b}\end{equation*}

这里为了便于分析，采用复数表达式.玩;'卢， $\omega$ 都是复数. c. c. 是式 (9. 14a) , (9.14b) 右边第一项的复共辄，因而 u ， ， $\rho$ 仍是实变量.由于方程 (9. 13 川， (9.13b) 是线性的，一般初值问题 (9.11) 可以用式 (9. 14a). (9. 14b) 形式解的线性 叠加来求得.因此只要分析 (9. 14 剖. (9. 14b) 形式解的性质，便可以得到稳定性的 判据. 将式 (9.14a) 和 (9.14b) 代入式 (9.13a) 和 (9.13b) 得到礼 ， ft 的方程如下 z

11 、 rJ U ， rJJj, 1 rJ 2iii -MEE ， +U， L二十 U ， 一一--二互+一一一一-\} iJx\} , -\} rJ x\} r7 x ,' Re r7 x\} rJ x\} (9. 15a) (9.15b) rJ ii; \_ iJ Xi 以及流动边界 Z 上的边值条件 z

\begin{equation*}iii(X ,t) = 0\tag{9.16}\end{equation*}

方程式 (9. 15 剖， (9.15b) 和 (9.16) 组成本征值问题，即在给定 Re 下.当 $\omega$ 取某 些本征值时，方程组才有非平凡解.本征值 $\omega$，的虚部可以用来判断流动是再稳定.将 复数 $\omega$ 写成

\[\omega =\omega ,  +i\omega ,  (9. 17)\]

式 (9.14) 中指数部分可写作 z exp(- i$\omega$ t) = cxp( 一 iw ， t )cxp$\omega$it (9.18) 因此若本征值的虚部 $\omega$i>O ，扰动将无限增长 z$\omega$i- 0 ，扰动是常幅振荡阳， <0 ，扰动 将渐近衰弱.如果在给定 Re 下一切本征值的虚部都满足 $\omega$i<O ，这时基本流是稳定 的，否则基本流是不稳定的.下面应用稳定性分析方法研究平面泊肃叶流动的稳 定性. 4. 平面泊肃叶流动的稳定性 平面泊肃叶流动的基本解是定常平行流，如图9.4 所示，如采用普通直角坐标系表示流场 2

\[U ,V , W = (U1 , Un U:i)\]

U(\_I') x \{x , y , z\} = \{XI , X2 ,X:I\} 图 9.4 平面泊扇叶流动的基本流 则基本流可表示为(见第 8 章):

\begin{gather*}
U = 1- y2\\
V=W=o
\end{gather*}

它的线性扰动方程可写作 z

\[iJ ii , , dU   - iJ jj , l' iJ2 ii , r1 z ii , rJ 2 ii  \]

一 iw ${\rm I}$i +U 一一+ ( --:$\tilde{li}$${\rm i}$ =- ':r +\textasciitilde{} (一寸 +-T 十.\textasciitilde{} : 1 r1 XI 飞 dy r rJ x ' Re 飞rJ x" r7 y" ' r1 z. I

\[eJ v rJ jj , 1 ,rJ 2 ii , r1 2 ij , r1 2 ii  \]

一 i$\omega$ 也 +U 一一=一旦旦+\textasciitilde{} (一$\tau$+ 一$\tau$+ 一$\tau$lrJ.l'1 iJ y , Re\textbackslash\{\} r1 x" ' r1 y" ' r1::;" J iJ jj , l' 俨毡:' I J2 U' I rJ 2 ii' \textbackslash\{\} 一 i$\omega$ ${\rm u}$J +U 一一=\_ ':r + 一卜丁+一$\tau$+ 一寸 ir1.1"1 r1 z ' Re 飞${\rm e}$J P 'r1 y" rJ z" J 空+豆豆 +$\tilde{w}$ = 0 ${\rm e}$J x ' r1y , ${\rm e}$J z (9. 19a) (9. 19b) (9.20a) (9.20b) (9.20 c) (9.20d) 线性方程组 (9.20) 的系数 U 和 dU/dy 都不含变量.r .z ， 因此该方程的解可以进一步 分离变量为

\section*{9.2  流动的稳定性分析}\addcontentsline{toc}{section}{9.2 流动的稳定性分析}

\{革，旬，命 .${\rm þ}$\} = 恼，台 .w. 如 exp(i$\alpha$x+i${\rm ß}$z) (9.2 1) 这里 ${\rm u}$( 川，台 (y) .${\rm U}$J(y) .${\rm þ}$(y) 只是 y 的函数，并称为扰动的模态或振型，假定 $\alpha$ .${\rm ß}$ 只 取实数，它们分别表示周期性扰动在 z 及 z 方向的波数.即相应的波长为 2$\pi$ .r- 一 ${\rm •}$ ${\rm A}$. = 一 (9.22) $\alpha$ 卢 进一步假定扰动是二维的，即 w三 O.${\rm ß}$=O. 这时扰动的振型须满足以下方程 z 一以+川 + (DU) 乡=一 iap+  (DE-d)\& 一 i$\omega$ -${\rm u}$+ 趴=一 DD+  (DZ-t i$\alpha$ ${\rm u}$+D 台 =0 式中 D=d/dy(D 是微分算符).扰动方程的边界条件为 (9.23a) (9. 23b) (9.23c) 元(士1)=识::!::1) =0 (9.24) 方程 (9.23).(9.24) 组成一本征值问题，它可以进一步简化为标准的常微分方程 的本征值问题.首先在方程 (9. 23 川， (9.23b) 中消去压强振幅户，然后利用连续方 程 (9.23c) 消去 L 得到如下的四阶复常微分方程 z $\downarrow$ (D2 一 $\alpha$ 2 ) 2 -${\rm u}$ + (U - c)( D2 - a2 )台一 (D2 U) 心 =0 (9.25) al( e 这里 c= $\omega$/$\alpha$ ，它的实部 c= $\omega$r/ $\alpha$ 表示扰动传播的相速度.式 (9. 25) 称为奥尔-索末菲 尔德( Orr-Sommerfeld) 方程，简称 O-S 方程. 利用连续方程 ia${\rm u}$+D -${\rm u}$ =O. 方程 (9.2 日的边界条件可由式 (9.24) 导出为 以::!::1) =D 以::!:: 1) = 0 (9.26) 对于方程 (9. 25) 和 (9.26) 组成的本征值问题，原则上可由如下步骤求得它们的解，四阶 常微分方程 (9.25) 应有四个线性独立的特解仇，它们都含有参量 $\alpha$ 巾 .Re. 即 ';i =骨i(y;a. $\omega$ .Re) ， i = 1.2 ,3.4 于是方程 (9.25) 的一般解为 (9.27) 台 = "!Ci';i (9.28) 式中 Cj 为任意复常数，将式 (9. 28) 代入边界条件 (9.26). 得以下的齐次线性代数方 程组 z C1';1 (1 ;a ， $\omega$ .Re) + C2 仇(1 ;a. $\omega$ .Re) + C3 仇(1 ;a. $\omega$ .Re) + ,A. (1 ;$\alpha$ ， $\omega$ .Re) = 0 C1 仇 (-1 ;$\alpha$ ， $\omega$ .Re) +C2 仇 (-1;a. $\omega$ .Re) 十 Cd'3 (- 1 ;a. $\omega$ .Re) + C4';4 ( 一 1;a. $\omega$ ${\rm •}$ Re) = 0 C1 趴(1 ;a ， $\omega$ ， Re) + C2 码(1 ;$\alpha$ ， $\omega$ .Re) +C3.;\textasciitilde{} (l ;a. $\omega$ ，Re) +c..;\textasciitilde{} (l; $\alpha$ 巾 .Re) = 0 C1 .;; (一 1;a.w.Re) +C2';\textasciitilde{}( 一 1;a ， $\omega$ ，Re)

+C3 钊 (-1;$\alpha$ ， $\omega$. Re) + ('1 1>\textasciitilde{} (- 1 ;$\alpha$ ， $\omega$ .Re) = 0 这一线性代数方程组有 Ci 非零解的条件是系数行列式等于零，即 1>1 ( 1 ;$\alpha$ ， $\omega$ .Re) 1>2 (1 ;$\alpha$ ， $\omega$ .Re) 仇 0;$\alpha$ ， $\omega$ ， Re) 仇 0;$\alpha$ ， $\omega$ .Re\} 1>1 (- 1;$\alpha$ ， $\omega$ .Re) 1> 2(-1;$\alpha$ ， $\omega$ .Re) 仇 (-1 ;$\alpha$ ， $\omega$ .Re) 仇( 一 1 .$\alpha$ ， $\omega$ .Re)

\[1=0\]

叫 0;$\alpha$ ， $\omega$ .Re\} 1>\textasciitilde{} ( 1 ;$\alpha$ ， $\omega$ ， Re) 1> \textasciitilde{}1 (1 ;$\alpha$ ， $\omega$ .Re) 1>\textasciitilde{} ( 1 ;$\alpha$ ， $\omega$ .Re) 到(一 1 ;$\alpha$ ， $\omega$ .Re) 到 (-1;$\alpha$.$\omega$ .Re) 1> \textasciitilde{}(-1;a. $\omega$ ， Re\} 1>: (一 1;a. $\omega$ .Re) (9.29) 等式 (9.29) 实际上表示 $\alpha$ ， $\omega$ .Re 之间的本征值关系式: $\phi$($\omega$ 心 .Re) = 0 (9.30a) 将此隐函数显化(或数值化〉后得

\begin{equation*}\omega =\omega (\alpha  ,  Re\tag{9.30b}\end{equation*}

它的虚部表征扰动的增长或衰减:

\begin{equation*}\omega =\omega i (\alpha  .Re)\tag{9.31}\end{equation*}

正值的创表示扰动增长，负值 $\omega$，则扰动衰减 .$\omega$，和 $\alpha$ 、 Re 有关，因此扰动增长与否依赖 于参量 Re 数和扰动的波长 ).=2$\pi$/ $\alpha$. 定义 9. 1 扰动的中性曲线 在 $\alpha$ .Re 平面上函数

\begin{equation*}\omega i (\alpha  .Re) = 0\tag{9.32}\end{equation*}

称为中性曲线.参数 ($\alpha$ .Re) 在中性曲线 一 侧 是增长扰动，另一侧是衰减扰动(见 图 9.5). 从图9.5 巾可以看到，当 Re 小于中性 曲线上的最小 Re 时，一切被数的扰动都是 衰减的.根据稳定性的定义，这时基本流在 线性条件下是稳定的. 定义 9.2 中性曲线上的最小 Re 为线 性临界 Re. 简称临界雷诺数.并用下标 cr 表 示: Re口$\cdot $ 稳 \textbackslash\{\} 定、 医 ( Rt') c 营

\begin{gather*}
----\omega 'i <()\\
R<!
\end{gather*}

图 9.5 平面泊扇叶流动的稳定性(示意) 虚线表示稳定 $\omega$， (a.Rp ) < lJ; 实线灰水 不稳 定阳i (a.R\textasciitilde{}) ? (); 垂 $\tilde{I}$ 线卖 iJ\textasciitilde{} 11 也界需谱数 这样就得到基本流的稳定性判据:当 Re 小于临界雷诺数时基本流是线性稳定 的.否则是线性不稳定的.对于平面泊肃叶流动，已由 ()- s 方程本征值问题的数值解 求出:

\begin{equation*}Re n = 5771\tag{9.33}\end{equation*}

就是说当 Re<5771 时， 一 切波数的小扰动都是衰减的 ;Re > 5771 时，在某 $\cdot $ 波数范 围内(巾性曲线内部)扰动是指数增长的，因而基本流是不稳定的. 利用线性稳定性还可以得到扰动发展的形态，当本征关系式 (9.30) 解嗣后，扰动

\section*{9.3  湍流的统计理论}\addcontentsline{toc}{section}{9.3 湍流的统计理论}

可表示为

\[u, = u ,  (y)exp[i\alpha ( .r - ct) J + c. c.\]

$\rho$= 户 (y)exp[ia( .r - C/) J + c. c. 其中 C= $\omega$/$\alpha$=$\omega$r/$\alpha$+i$\omega$j$\alpha$ ，也可将式 (9.34 时，式 (9.34b) 写成

\[u, = u, (y)exp[i\alpha ( .r - Crl) Jexp\omega ,  1 + c. c.\]

(9.34a) (9.34b) (9.35a) $\rho$= 卢 (y)cxp[i$\alpha$(x-cr/)Jexp$\omega$ ， 1 + c. c. (9.35b) 除了增长(或衰减)因子外，扰动以相速度 : C r = $\omega$r/ $\alpha$ 传播.在平行基本流中 $\omega$r 手 O. 而 且已经有人证明 O<Cr<(U(y) )m. 萃，因此扰动是以行进披方式传播，它的传播方向和 基本流方向相同.在平行流 ， t 1 线性扰动的行进设又称托尔明-希利希延 ( Tollmien Schlichting) 波(简称 T-S 波) . 本节的方法可以应用于其他平面平行基本流的线惶稳定性分析，对于不同速度 分布的基本流，它们的临界参数 Re${\rm ‘}$ v 及特征状态(如 $\iota$F 等)是不同的，但扰动以 T-S 波 的形式传播是共同的.同时这一分析方法也可以近似地应用于基本流是缓变的准平 行流. (U(口，川 .EV(U.y) .O\}. 例如层流边界民流动等. 最后，应当注意，线性稳定性分析的结果是在微弱初始扰动条件下考察基本流的 稳定性.严格地说是在..零"扰动条件下的稳定性判据.因此它总是大于实际上开始发 生不袁减扰动的 Re. 计入非线性效应后 .1临界 Re 将降低并接近实际观测值.另外当 初始扰动增长以后，非线性效应很快起主宰作用，它将导致 - 维的、多颇段的扰动直 至发展到有宽带连续谱的充分发展湍流.由于目前缺乏有效的数学分析 t 具来预测 这种复杂过程.关于流动由 Jzf 流过$\cdot $渡到湍流的全过程尚未完全了解.读者可以在流动 稳定性的专著或期刊论文小学习这方面的内容.

现在来讨论充分发展的湍流 .9. 1 节 '1' 已好-说过.湍流的主要特征是含有宽带连 续谱的不现则流动.对于不现则现象.人们通常应用概率分布的方法研究它们的统计 规律性. \{E 纯典 1由流理论 '1' ，需诺 (1894) 芮创J-I J 时间平均的方法把不规则部分..过滤" 掉，然后分别研究平均部分(确定性的)和脉动部分C:\{\textasciitilde{}现则的 h 并在此棋础上建立预 测湍流平均运动的统计理论.随着湍流研究的深入.人们发现时间平均方法有局限 性，官只适 JH 于平稳的湍流( 1(: 时间平均有l 起始时刻无关) ; 例如.湍流边界层的来流 是周期性的流动属于 11: 平稳油流.对于这种非平稳揣i 流，用 K 时间平均方法.将使周 期性变化初不规则湍流 一 起过滤 J申'.现代湍流应用系综平均概念.将一次湍流试验的 流场作为一个样本.在相同的边界条件 F. 重复无数次试验的样水平均称为系综平

均，在实际物理实验或数值模拟过程中，以足够多的样本平均作为系综平均.下面介 绍统计平均方法，然后导出统计平均方程. 1.瑞流的统计和分解 定义 9.3 流动参量的统计平均值等于它在给定边界条件下一切可能出现的流 动事件的平均值. 其数学表达式为 (U;) = 为主 uj3 (9.36a) N (P) = 古 Zpfn3(93仙 Uj 川， p C ${\rm •}$ ，为给定边界条件下一切可能出现的流场速度和压强 ， N 是统计的样本数.系 综中任意事件的流动参量和它的平均值之差称为脉动量，其数学表达式为 以 =Ui 一 (Ui ) (9.37a) 户'= P 一 (P) (9.37b) 如果系综平均和时间无关，在统计理论中称之为平稳态，在流体力学中称为统计 定常的，简称定常湍流.统计定常的湍流可以用长时间平均来代替系综平均 z (U,) = !\textasciitilde{}\textasciitilde{}[ $\tilde{f}$' Ui( 川世] (9.38) 时间平均方法可以看作一种过掠过程.图 9.6 形象地演示平均过滤. 由平均值定义式 (9.36) ，式 (9. 37) 很容易导 u 出以下平均量的性质. (1)平均值的系综平均等于平均值本身.

\begin{equation*}((U ,)) = (U;)\tag{9.39}\end{equation*}

(2) 脉动值的平均值等于零.

\begin{equation*}(u) = 0\tag{9.40}\end{equation*}

式 (9.39) 直接由式 (9. 36) 导出，因为系综平均时 已经包括了一切可能发生的状态，对它再取平均

\[r+T '1'\]

图 9.6 雷诺平均的示意图 必然等于自己.式 (9.40) 由式 (9.37a) 和式 (9.39) 导出，将式 (9. 37a) 两边取平均

\[(U i ) = ((U;)) + (u:)\]

利用式 (9.39) ，便得式 (9.40). (3) 脉动量的一次式与任何平均量的乘积的平均值为零，即

\begin{equation*}((Q)(P) \rho ')=0\tag{9.41}\end{equation*}

这一等式也很容易由定义来证明(读者自己证明).但是一般来说脉动量的 n 次乘积 (n> l)的平均值不等于零，即: (u;u; > , (u' p') , (u;u; … u\textasciitilde{}) 等一般并不等于零.

(4) 平均运算和导数运算可以交换 〈旁> =手. <窍> =譬 (9.42) 因为求导运算与求和运算可以交换，直接由平均的定义可以证明式 (9. 42) 如下: 〈苦> =走去 22::= 去(为主 QC"))= 孚 〈主> =力去哇!:=(为主 QCn) )=守 以上诸性质在建立湍流统计理论时将经常用到. 2. 脉动量的关联 脉动量的平均值等于零.因此平均值不能反映脉动量统计性质的差别.为了考察 脉动量的统计特征和不同脉动量之间的统计关系，利用脉动量乘积的平均值.并称之 为关联. 定义 9.4 同一脉动量 n 次乘积的平均值称为 n 阶自关联;不同脉动量的 n 次 乘积的平均值称为 n 阶豆关联. 关联是脉动量之间统计上的联系程度，如两个脉动量的关联等于零，则称它们在 统计上不相关.或统计独立.下面我们介绍几个常用的关联. (1)揣动能 定义 9.5 脉动速度平方的统计量之半定义为流体质点单位质量的揣动能 k. 简 称揣动能 z h=$\div$ 〈以〉 (9.43) 揣动能是脉动连度的自关联.很容易证明.流体质点单位质量的动能平均值等于平均运 动的功能和揣动能之和.令 K=Ui U，I 2 表示流体质点单位质量的动能，则其平均值为

\[(K) = ~ <U,U,) = ~ ((Ui)<U ,) +2u:(U i ) +u:u) 2\]

利用前面性质 (9. 39) 和性质 (9.4 1).得 〈 K 〉=$\div$〈 U ，〉〈 UJ+$\div$ 〈 umb(9$\omega$ 式巾 (U ， )<U ， )/2 是流体质点单位质量的平均运动的动能，因此揣动能 h 把脉动具有 的动能从质点的平均动能中分离出来. (2) 湍流度 e 定义 9.6 脉动速度的均方根与当地平均速度绝对值之比称为湍流度，并用 r 表示 z 2-| 、

\[3r-\]

m 、、/ H 川-E u 一

\[-II\]

1ll

\begin{equation*}r\tag{9.45}\end{equation*}

式中 (u;u) = ( u; u; +u$\tilde{u}$; + u$\tilde{u}$\textasciitilde{}). 湍流度表示当地脉动速度的相对强度.在实验流 体力学J.j l ，还常用 e=(u;u; /3 )1 2/I( U ， )1 表示湍流度.它的含义和式 (9. 45) 没有根本 差别. (3) 关联函数 我们还可以利用关联来考察脉动量在时间序列上或\textasciitilde{}间分布的统计特征.例如: ${\rm 1}$ 二 阶时间关联. 定义 9.7 同一宅间位置上，相隔给定时间差 r 的两个脉动 ht 时间序列之积的 平均值称为这两脉动量的二阶时间关联，用 Rj (X ,l ;r) 表示，它的数学表达式为 R '1 ( X ,1 ; r) = ( u: ( x ${\rm •}$ t ) u: (X.l + r) ) ( 9. 46) 这 一 关联量表征两脉动量在时间序列的发展过程中的统计关系.如果 Rj (r) = 0 则 表示 4.u: 这两个脉动量在相隔时间 r 上互不相关.同 一 时刻的统计量自关联总是大 于零，例如，揣动能，它是自关联 R" (x , J) = (u; (x. t) u; (x. 1) ) 之半.总是正值. 一 般来 说.两个随机函数在时间上相隔很长以后就互不相关，因此通常有 R" ( $\infty$ )三 O. ${\rm 2}$二阶专间关联. 定义 9.8 同一时刻，相隔给定空间位移 r 的两脉动量之积的平均值.称为两脉 动量之间的 二 阶宅间关联.记作R" (X.I ; r) ${\rm •}$ 即 $\tilde{N}$ R" (x.t;r) = ( U:(X. J) l/(X 十 r.t)) (9.47) 这 $\rightarrow$ 关联量表征两脉动量空间上的统计关系，女n 果 R ij (r) = O. 则表示 uf ， u; 这两个脉 动握在雪间方向 r 上互不相关. 一 般来说，在空间上相隔很垣的两个随机量旦不相 关.因此有 R$\upsilon$( $\infty$ )三 O. ${\rm 3}$二阶反 n 阶时空关联. 由前面定义的时间关联和空间关联，不难把关联的统计推，..到四维时专以及把 关联量增加到 n 阶.本书不打算详细讨论它们 .1对为 - 般仁程庇用'1 1 常常遇到的是 二 阶关联，至多 三 阶关联，例如， 二 阶时专关联: R$\upsilon$ (x ， l;r ， r) = (u:(x.t)u:(x+r ， l 十 r)) (9.48 (1) 又如， 三 阶空间关联:

\begin{equation*}R,j.k(x ,t;r) = (U:(X ,  l)U:(X.l)U~(X + r.l))\tag{9.48b}\end{equation*}

在 三 阶相关巾，为了表示脉动量在不同的时空点上，常用远号将不同点区别开来，自日 式 (9.48b) 叶'u; 与 uJ.u; 在不同的空间位置上. 3. 统计平均纳维-斯托克斯方程及脉动方程 湍流是不规则运动，宫的脉动部分是随机性的，但是统计平均量是规则的，而人 们通常需要知道 t 程和自然环境中湍流运动的平均特性，有了统计平均方法，可以对 湍流场中平均量进行定量的研究.具体做法是对不可压缩牛顿型流体的运动方程逐 项做系综平均:

〈等) + (U， 去) =一〈去〉十趴在;〉 〈在) = 0 由于系综平均和求导是可交换的，因此平均运动方程中的线性项可利用式 (9. 42) 简 化，平均的连续方程可简化为平均运动的连续性方程: cJ (U,) \textasciicircum{} rJ Xi 平均非定常项、平均压强梯度项和平均粘性项可分别等于平均量的导数.即 〈等) =嗖二 〈旦\textasciitilde{})一旦旦-

\[aXi \theta  Xi\]

J 一孔，

\[U-J a-h\]

一一 、飞/'U-h 俨一几，，，、、 为了导出 (UjdU;/ .J:r，) ， 利用连续性方程U;/${\rm U}$ :r i=O. 可得 (U 旦旦\textbackslash\{\} = I 旦旦 -U ， $\tilde{U}$， '> = (町i U , rJ.l' j I 飞 ${\rm U}$Xj -, rJ.1"j I \textbackslash\{\} d.1."j =是 ((U， )川 + u\textasciitilde{} (Uj ) 十川;十以〉 利用平均值性质式 (9. 39)\textasciitilde{} 式 (9.41) ，上式可简化为 〈 UJ?$\upsilon$) r .-,. ' -\}' . rl U i ) = 乒 (U ， )(Uj ) + ::J'\textasciitilde{} (川 j) 时，l'j , rJI\} rJX) J(U ,) 利用平均运动的连续性方程一一':":"'=0 ，上式右端第一项可简化为rJ.1.", 主 (U ， )(U ，) = (U，) 主 (U;)

\[<l.1."j "'.1.",\]

将以上关系式代入平均运动方程后得

\[J(U;) I JTT   rJ (U ,) '"-- ,J( P) I U\]

2 (Ui ) ${\rm e}$J -一$\iota$ + (U ; ) 一一....!....:.\_ :i;::一一一二十 $\nu$ 一一-一一千一叫 u\textasciitilde{})

\[J' u::r:, e).1',' - dXj rJ Xj C) Xj . -- ,--)\]

连同平均运动的连续方程 (9.49a) rJ (U;) \textasciicircum{} ${\rm e}$) Xi (9.49b) 构成平均的纳维-斯托克斯方程，也称雷诺方程，因为雷诺最早应用时间平均方法导 出该统计方程.如果我们把平均运算施加于用应力表示的不可压缩流体的动力学方 程，则雷诺方程可以写作 z rJ (U;) I JJT\textbackslash\{\}(J( U ,) \_ J(P) I ${\rm u}$ /町'..' \textbackslash\{\} 一一$\iota$ + (U ，) 一一一一一-一一+一一〈一 - u';u' , ) (9.50) Jt .. - J ${\rm •}$ dX, ${\rm u}$x,' ${\rm e}$J .T， 飞 p --,-- J I

式 (9. 50) 不限于牛顿流体.也可应用于非牛顿流体.式 (9. 49) 和式 (9. 50) 表明.雷诺 方程在形式上和纳维-斯托克斯方程极其相似，只是对于平均湍流场来说应当在平 均应力项加上一项 -p 〈 uh; 〉，该项称为雷诺应力，我们将在下面仔细研究它. 对于平均运动方程来说雷诺应力项是未知项，因此雷诺方梧的未知量有 (U ，) ${\rm •}$ (P) 和 (U;U; )共 10 个.超过平均运动的方程个数，因而雷诺方程是不封闭的.只有补 充了雷诺应力的方程后.才可能使雷诺方程封闭，并用它预测平均场的运动.这是揣 流统计理论需要解决的主要问题 雷诺应力的封闭方和. 将原始的纳维-斯托克斯方程和雷诺方程相减，可以得到脉动场方程如下 z 'J,< UF >二子L 十 (U. ) 二三\textasciitilde{}+U' "，，，一- dl dIJ dIJ ab' , 1 世'=一二.L+一一旦j\_ - \_j!\_( ($\mu$';) - U 川) a .1::i . Re a.l'J a.l"J c7 x J ' ${\rm •}$ -- ,-- J ${\rm •}$ -- -- J (9.51a) iJ u$\tilde{n}$

\begin{equation*}<7.r,\tag{9.51b}\end{equation*}

请读者自己导出此公式.同雷诺方程一样，由于存在雷诺应力 (U:U;) ${\rm •}$ 脉动方程也是 不封闭的. 4. 雷诺应力 上面导出的雷诺方程中有一项不封闭量 (U;U;). 通常称 -p(u:U;) 为雷诺应力，为 书写简单起见，以下常用 R'J 表示雷诺应力.在不可压缩湍流中 p 是常数，有时省略密 度 p. 而直接将一 (U:U; )称作雷诺应力，即R.j = - (U;U;>. 需诺应力有以 F 基本性质. (1)雷诺应力是二阶对称张挝 (U;U; )是一点(即 r=O.r=O) 的二阶互关联.它有乘积可交换性 : (U:U;) = (U;U 刀， 即有

\[R'J = Rj,\]

这就是说雷诺应力和分子粘性应力相仿是二阶对 称张量.它有六个独立分量. (2) 雷诺应力是单位面积上脉动动量通量平 均的负值 我们取一微元立方控制体(图 9. 7). 通过法向 量为 i 的控制面上单位面积的动量通量为 Mi = 乒JiU 它的平均值为 (9.52) 4乙\_!\_\_\_\_/ I v, = W, ) +屿' 叫乡，丁'"7 $\iota$ = (U)情' |每1 9.7 说明脉动 4IJ 量用|| U 是速度向量 ， V ， 是 z 方，:.，的速度分屉

\[<M,> =p<(<U,> +u:)( <U >+uF) >\]

利用平均值关系式 ((Ui)U;)=O. 平均动量通量

\begin{equation*}<Ma > =p<U,>< U > +p<u(u' >\tag{9.53}\end{equation*}

式 (9. 53) 右边两项分别表示通过控制面的平均运动的动量通量 $\rho$ (U;) <U) 和通过控 制面的脉动动量通量的平均值 p(u:u') ， 因此式 (9.53) 表明: 通过控制面单位面积上的动量通量的平均值 =通过该控制面的平均运动动量通量+脉动动量通量的平均值 如果我们在控制体上建立平均运动的动量方程，上面的等式说明对于平均运动 来说，在控制体表面上除了有平均运动的动量通量外，还有平均的脉动动量通量，根 据动量定理，相当于在控制界面上有 一 当量的表面力强度，它的数值等于脉动动量通 量平均值的负值，它就是雷诺方程中的雷诺应力. 5. 雷诺应力输运方程和瑞动能输远方程 雷诺应力既然是脉动动量通量的平均值，它必然遵守一定的动力学关系.事实 上，我们可以由脉动方程导出雷诺应力方程.具体步骤是:用 u; 乘 u: 的动量方程加上 用 u: 乘 u; 的动量方程，再取平均后经适当的代数推导可得如下雷诺应力输运方程 z t 一丁7i$\iota$一 + ( U叭k) 一7一L - 一 -<k r  :u;> EjEZ <U:U \textasciitilde{})守 + 旦去且三 :\} I (P $\gamma$ ) " ， (p $\gamma$).. I 1..'..'\_ ,' \textbackslash\{\} J(u$\tilde{u}$:) 飞 / Ju$\tilde{Ju}$: \textbackslash\{\} 一 $\tilde{I}$ 一一$\rightarrow$J 品+一:...:::..L:...${\rm O}$ ${\rm u}$< 十 (u:u:u $\rho$ 一 $\nu$ 一一一L:... J-2$\nu$ 〈一一〉r1.J， 'k 飞 p -,. , p -\textasciitilde{}， '--'--'-'" ${\rm •}$ J.rk I 飞 J .r k J.rk / E '1 (9.54) 将雷诺应力输运方程作张量收缩，因有 (U:U:) =2 走，可得到揣动能方程如下 z r1 k , ,. T , Jk , , " J(U;) 一 + (U.) 一一 = 一 (u:u: )一一 Jt ${\rm ‘}$ $\theta$r品，--. ' J.rk 「矗 盛 , 1 $\theta$ I ($\rho$ 'U\textasciitilde{} )$\theta$k \textbackslash\{\} / Ju$\tilde{Ju}$\textasciitilde{} \textbackslash\{\} 一一:.\_r 一-一一+怡 'u:)-IJ\textasciitilde{}一-\}-IJ( 一一一一) (9.55) $\theta$Z矗飞 p … J.rkl 飞 J .r k $\theta$ .r k / D h E 方程 (9. 54) 和方程 (9.55) 分别是雷诺应力和揣动能在平均流场中的输运方程.方程 中各项的物理意义可简要地解释如下. (1)两方程左端是沿平均运动轨迹雷诺应力或揣动能的质点导数，也就是沿平 均运动轨迹上雷诺应力或揣动能的增长率，也就分别用 C;j 和 C k 表示.

(2) 两方程右端第 一 项 P ;j 和 P. 称为生成项，它表水平均切变场与雷诺应力相 互作用对于雷诺应力或揣动能增长率的贡献. (3) 式 (9. 54) 右端第 二 项矶是压力脉动与脉动切变率的关联，称为再分配项. 因为在不可压缩流动中.脉动速度的连续方程为 c7 u: / r7.r ， = o( 见式 (9.51b)). 故 $\phi$'J 作 张量收缩后等于零，即也 =2($\rho$ 'au; / J :J.';) =0. 也就是说，再分配项矶对揣动能的增 长率没有贡献，它的作用只是在雷诺应力的各个分量 之间起调节作用.下面将具体说 明 $\phi$$\eta$ 的作用. (4) 雷诺应力或揣动能输运方程右端第 三 项 D'J!或 D. 是散度形式，称为扩散项. (5 )雷诺应力或揣动能输运方程最后 一 项是起散项，因为在揣动能方程叶 1 Ek = 一 $\nu$ 〈 $\theta$ U; / c7.r k c7 u : / .J.r$\rho$ < 0 ，它总是负值，即这 \_ .项总是使揣动能衰减.揣动能衰减常用 E 表示， $\epsilon$ = $\nu$ ( r1 u: 闭 .Tk r1 u: / r7 .rk ) . 研究雷诺应力或揣动能输运方程的各项平衡关系是湍流理论的专门内容.下面 简要地分析 一 下揣动能方程的平衡关系. j由动能的主要生成项是 P. = - ( u; Ii \textasciitilde{} ) r1 ( U, ) / r7 .1". ，它是需诺应力在平均切变场 上所做的功，如果平均速度场是均匀场，即 以 U ， > = 0 那么就没有揣动能生成 (P k =0). fflii揣能总是耗散的，即 - E < o.1大!此在均匀的平均速 度场中揣动能将 一 直衰减.换句话说，必须通过平均切变率才能由平均场向湍流脉动 输送能址. 在 j白流统计理论中.揣能起散 卫字， E: 是各种脉动成分艳散的平均值.在牛顿型流体 中揣动能在散率为 E = 川 r1 u; / rJ.rkc7 u: / rJ.1". > .脉动速度梯度 I rJ U: / rJ..l" kl 越大，对平均揣能 起散率的贡献越多.从量级 11${\rm i}$ ${\rm i}$ ' i-有 I r1 1/!rJ.l"k I \textasciitilde{}!k/ 1. 其'1 1 、!k 是 脉动速度的 量 级 .1 是 脉动的唁间尺度.由以上的 it 级估计可见.小尺度的脉 4J 成分对揣能耗散起 主 要作 用，或者说.大部分揣动能在小尺度脉动巾槌散.于是从揣能生!现率f.IJ 桂散率的简要 分析中 I可以引什l 以下的概念: i市流脉动从平均切变场(大尺度)申得到能量而在小尺度脉动中耗散. 在定常的缓变平均流'1 1 .m. 平均运动轨迹揣功能的变化率以及扩散率部很小而可 略去不计.剧l 假定 C.Dk 近似等于 零 .我们称这种揣 i成为平衡流动.于 是 在平衡流 '1 1 有 端动能的生成率 1飞 = 端动能的耗散率$\epsilon$ 具有平衡流性质的湍流 是 比较简单的湍流流动.称为简单切变湍流. 对揣动能的产生有了基本的认识后 .11 j.简怪地考察帘谣i J缸力的输运过程.雷诺应 力的扩散和在散的过程与揣功能在输运过程巾的作用相仿.着重探讨同分配 r${\rm o}$i 钱， 的 作用.为了简单起见.考察 二 维的均匀简单剪切流 z 平均速度场 < U ， > =U( 川 ) ${\rm o}$'d ，同

时rJ ( ${\rm •}$ )/rJ XI = 0 和rJ( ${\rm •}$ ) / rJ X:I =O. 以及 ( U: U: ) = (U$\tilde{U}$: ) = O. 困此.只要研究 (U; U;) 、 (U$\tilde{U}$ \textasciitilde{} )和 (U;U\textasciitilde{}> 三 个雷诺应力分量的输运方程就够了.将二维平均场公式代入雷诺 应力输运方程 (9.54) .就有 .1 (u$\tilde{u}$\textasciitilde{} 飞 1 ， dU(x ,, ) 一一-一一=- 2<u$\tilde{U}$\textasciitilde{} > 一气--一 +$\phi$11 +D"-E ,, .11 -.--,--\_ . dx . -"… E (9.56a) a < u$\tilde{u}$ 门" dU(.r z ) 一寸\textasciitilde{} = \_ (u$\tilde{u}$ ;) U'-'I" 一 +$\phi$1 2 + DI2 - EI 2 dl dX 2 (9.56b) E D + $\phi$ + nu --川7 (9.56c) 由式 (9.56a) 可以看到如果存在雷诺剪应力 - (u:u\textasciitilde{} )> O. 则通过平均剪切率场的作 用 (P 1 2 = - 2 ( u; u\textasciitilde{} >dU / d.l. 2 > 0). 流向的脉动速度将增1.(;而由式 (9.56b) 可以发现， 要使雷诺剪应力增辰，必须存在横向的脉动 ( u $\tilde{u}$ ;) ; 但是式 (9.56c) 告诉我们 二 维平 均剪切率场对 ( u$\tilde{u}$; )没有贡献.假如没有产生横向脉动机制，那么 ( u;u;) 就不能增 长.也就不可能维持雷诺应力 -(u:u \textasciitilde{} ). 上面的分析说明，仅有平均剪切事和雷诺应 力生成项 P ij 不足以维持需诺应力的输运.雷诺应力的再分配项趴在雷诺应力输运 过程 rll 起重要作用.当流向的脉动 <u; u;) 产生并增民后，因湍流脉动间的相互作用发 生脉动能量的再分配 $\phi$')? .它将 ( U'I U'I ) 中 一 部分动能传递给横向脉动 ( u 二 u 飞 〉 \{式 (9.56c)); 然后通过方程 (9.56b) 维持雷诺剪应力. 7，由流脉动间的能量传递是再分 配项 $\phi$'J 的主要功能.它是在湍流脉动分量间起调剂的作用，虽然'仨并不对揣动能的 增民有任何贡献. 总之，在平均剪切流'1 1 维持湍流的统计机制可概括如下:通过平均剪切率产生 流向脉动，通过再分配项将流向脉动的动能传递给横向脉动.横向脉动在剪切场中维 持雷诺应力的输运. 应当注意到.虽然需诺应力输运方程补充了六个雷诺应力的方悍，但是这组方程 本身义引进了新的高阶统计 -\&t: ($\omega$tut:〉u:〉u;$\upsilon$) ${\rm •}$ (协$\rho$ 'U:> ${\rm •}$ (怡$\rho$ ， ($\upsilon$rJ U叫4〈:$\upsilon$/川rJ X盯j + rJ U叫; /川J 玩)门〉以及 〈$\omega$CJ句Ju叫; / rJ .T飞$\eta$J'${\rm I}$t 更多的高阶统计哇.因此 .JH 统计方法导出的方程组永远是不封闭的，越是高阶的统 计方程含有的术知项越多.同此应 JH 统计理论来预测平均场时，必须对未知的统计量 作合理的假设.这比假设常常称为片闭模式，我们将在 9.4 节介绍 一 些常用的模式， 有关雷诺应力输运方程 H 闭模式的详细讨论是湍流的专门课程的内容. 6. 各向同性瑞流的瑞动能传输的串级理论和科尔莫戈罗夫 (Kolmogorov) 能谱 (1)均匀各向同性湍流及其能谱 均匀各向同性湍流是 一 种最简单的湍流.均匀湍流的含义是 z 所有ii白流统计量的

空间导数等于零. 根据揣动能输运方程，均匀湍流没有对流、扩散和生成，因此均匀湍流必定是衰 减的. 各向同性湍流的含义是:在任意转动的刚性坐标系中所有湍流统计量相等.也 可以说，在任意转动的刚体坐标系中观察湍流，所有的湍流统计址 是相同的. 关于各向同性湍流理论的严格定义和论述需要应用张 量 理论，不在本书范围， 读者可以参阅作者的湍流专著 《 湍流理论和模拟 >> .下面概要叙述各向同性湍流的 性质. ${\rm 1}$均匀各向同性湍流的衰减 均匀各向同性湍流中平均速度和揣动能满足以下方程:

\[eJ (U i > ~\]

${\rm e}$J X j (9.57a) (u:u:>" 2k ( u: 叫>=一一一」仇 = 一仇 (9.57b) J. 3 -'J 3 -'J 就是说，均匀湍流的平均速度场是均匀平行流，可以在以均匀速度移动的惯性坐标系 中观察均匀湍流，因为它们的动力学行为是相同的.均匀湍流场巾可以令 (U i > =0. 式 (9.57b) 表示各向同性湍流的揣动能均匀分布在 三 个方向的速度分量上，它的雷 诺应力等于零.将均匀湍流的关系式代人到揣动能方程 (9.5 日，可得 兰=一 $\nu$ 〈拦去) =-E 在前面已经论述过，不同尺度的湍流脉动它的揣动能的耗散率是不同的，小尺度的揣 流耗散率远远大于大尺度的湍流脉动.因此揣能的耗散过程中将不断由大尺度脉动 向小尺度脉动传输动能. ${\rm 2}$各向同性湍流的能谱 为了表示各种不同尺度湍流中所含揣动能的多少，引人能谱的概念.因为揣动能 是有限值，也就是说湍流脉动 u: 是平方可积的，因此它可以展开为傅里叶积分(理论 上来说， u; 是不规则的随机函数，它的傅里叶积分应是广义积分，把它当作常义积分， 并不影响我们对它物理意义的理解) : u: (.rl ,X 2 ,X 3 ' 0 = .ill ${\rm U}$i ( 川p( ik .川 = .ill ${\rm u}$,( 札川.k走乌川2$\upsilon$.k仇$\iota$们川3口川${\rm •}$ t川t${\rm ο}$\} 忧叫叫p证[ i(k丸忖k1xI卢山I盯I+k$\iota$2山I盯川2 + $\iota$h川r盯川! >汀Jdk$\omega$Id川3 式中品， (儿，鸟，儿)称作脉动速度 U ， 的波数谱，或简称谱肉，鸟，也分别为 .1' 1 ${\rm •}$ 岛， 町 方 向的波数.利用谱分析，可将脉动中不同尺度的成分用波数分解出来.小的披数相当 于挺披，因此小波数代表大尺度脉动;而大波数对应小尺度脉动.在谱分析中，物理量 在町，岛，町坐标系中的分布称为物理空间的分布，物理量在波数空间中的分解称为

谱空间的分布.由和、分表达式，显而易见 ${\rm u}$ ， (丸，走 2 ' k.. )表示在谱 空 间的微元体积 dk1dk l d 儿中脉动的强度.应用捕流统计撞在物理空间的各向同性的性质，可以证明， 湍流脉动在谱空间中统计量也具有各向同性的性质，例如，揣动能在谱宅间各个分量 上的分布是相同的.将它作统计平均，可得(证明从略，可参见湍流专著\}

\[E <t) = <U:U~) = fE \omega  .t)dk\]

式叶 1 k = v' k ; +k \textasciitilde{} + 肘，是波数向量的模，它只表示波数的大小，而不反映波数的方 向 .E(k.t) 表示揣动能在植数大小上的分布.称为揣动能的功率谱，简称能谱.注意 功率谱的量纲是单位质点动能乘以氏度.即 [E( 走 )] = C' T 2 波数表示不同脉动成分 的尺度，闲此能谱表示不同尺度的脉动在揣动能巾的贡献. (2 )各向同性湍流中揣动能传输的串级理论和能谱 ${\rm 1}$ 耗散尺度 前面已经说明，各向同性湍流中揣动能是衰减的.但是小尺度湍流脉动的衰减远 远大于大尺度湍流脉动的衰减，也就是说，能谱 E(k ， t) 的衰减率随波数的增大而急 剧增加.另 一 方面，低披数(大尺度)湍流脉动含有较大的能量，在揣动能衰减过程中 揣动能从低波数的成分向高波数成分传输，最后在很高的波数脉动中耗散.数学家科 尔莫戈罗夫首先用量纲分析方法预言各向同性湍流能 i普应有的规律. 科尔莫戈罗夫理论的前提是:在各向同性湍流脉动具有很宽波带，足以区分非 能散的波数带 k$\tilde{km}$"， 和挺散的波数带走 $\tilde{kJ}$' 即 kJ >> k m川设各向同性湍流的区域尺 度为 1..则最小的波数是 k nun = 2$\pi$ I L. 假定产生揣动能耗散的主要脉动尺度为守，则耗 散性的最小波数是 k J --1 1 r;.科尔莫戈罗夫理论要求 L >> l$\nu$/ 走k J ， 或 kJ >> k n川3 下面用量纲分析方法导出耗散区 1揣甜流脉动的量级关系.决定能散过程的物理量 是:湍流脉动的速度尺度 v( 可以采用揣动能的开方 .j2k) ${\rm •}$ 流体的粘度 $\nu$ 和揣能起散 率 E. 它们的量纲分别是:

\[[v] = MS 1, [\nu ] = M 2 S-I. [E] = M 2 S ~\]

由量纲关系 .v 、 $\nu$ 、 E 之间应有以下关系:

\begin{equation*}v - (\nu \epsilon )1"\tag{9.58}\end{equation*}

同样，由量纲分析可得起散脉动的尺度 l l k J - 市和耗散参数有以下关系 z 平 \textasciitilde{}($\nu$:' 1$\epsilon$)11 (9.59) 定义耗散的局部雷诺数，简称耗散雷诺数 (Re)D=V r; 1 $\nu$ ，则由式 (9.58) 和式 (9. 59) 关 系式可导出

\[(Re)v = vr; 1 \nu -1\]

就是说.能散雷诺数的量级等于l.在式 (9.58) 和式 (9.59) 中消去运动粘性系数 $\nu$ ，我 们可以得到揣动能耗散的如下关系式 z

\begin{equation*}E-v3 / r;\tag{9.60}\end{equation*}

根据科尔莫戈罗夫理论要求 L >> l / k J ${\rm •}$ 因此非耗散性的大尺度雷诺数 (Re) $\iota$=

\[vL/ \nu >> 1.\]

${\rm 2}$科尔莫戈罗夫揣动能串级理论和一 5/3 谱 科尔莫戈罗夫揣动能串级理论的核心思想是:假定前面论述的在耗散区揣动能 的传输过程具有相似性.就是说，从尺度 1/k 到尺度 1/(k 十 dk) 的揣动能传输都等于 揣动能耗散率 E. 并在 k ，，> 走 〉 儿'"的范围内能谱具有相似分布，即 E(k) = 旷/ 、 1 / 1 f( 是守) (9.61) 式巾 ，， ; / 1 E 1 川 是 E( 灿的量纲 .I( 忖)是无量纲相似函数.由于揣动能的耗散是从大尺度 脉动传输过来，在 k ，， >>kmlll 条件下，大尺度脉动(或低波数)传输过程应和流体的粘性 无关，因此，根据量纲分析，函数 I< 问〉必须是如下的幕踊数形式，可用量纲分析和 式 (9.59) 导出: f( 句 )=$\alpha$a(k 句守$\rho$) ，， 1 目:i =$\alpha$$\omega$$\nu$ 5川/仆1$\epsilon$$\nu$i.;;川12切k :, 代入式 ($\omega$9.61) .得能谱公式:

\[E(k) = \alpha \epsilon  2 / 3 k , / :<\]

式 (9.62) 是著名的科尔莫戈罗夫 -5 / 3 律.

\[< 9. 62)\]

实际中存在近似的均匀各向同性湍流，例如在风洞中均匀流动通过细网后的揣 流，以及开阔的大气、海洋中的湍流等;即使是不均匀的切变湍流，在局部小范围内， 湍流脉动也可近似为各向同性湍流.所有这些例子都可归结为局部各向同性湍流.在 上述流场中实验测量都证实存在一 5/3 律的湍流能谱.

\section*{9.4  湍流封闭模式}\addcontentsline{toc}{section}{9.4 湍流封闭模式}

\section*{9.3  节导出的湍流的统计输运方程中总是含有高阶统计矩形式的不封闭项.例}\addcontentsline{toc}{section}{9.3 节导出的湍流的统计输运方程中总是含有高阶统计矩形式的不封闭项.例}

如，雷诺方程是平均速度场的输运方程，它含有二阶统计矩不封闭量:雷诺应力 - (u;u;). 在雷诺应力输运方程 '1 1 则含有待封闭的三阶矩(例如 : ${\rm i}$由流扩散 D i\} )和更 高阶矩(例如 z 再分配项 $\phi$，JM]). 湍流模式的基本思想是在低阶统计矩的输运方程中 建立高阶统计矩和低阶统计姐间的关系式，一旦这一关系式建立以后，低阶统计矩输 运方程就封闭了.例如，在雷诺方程中建立雷诺应力- (u;u\textasciitilde{}) 和平均运动场 (Ui ) 间的 关系式，需诺方程就封闭了，这种模式称为低阶短模式.如果以雷诺方程和需诺应力 输运方程联立作为湍流统计平衡的控制方程，那么建立所有待封闭的高阶矩项和雷 [注] 可以证明.脉动压强含有 三 阶速度相关 J$\tilde{i}$.因此同分配项 $\phi$;j =( 户 '[(，) u: /.J .rj )+( <1 u;1 归 i) J) 含有四 阶统计矩.

\section*{9.4  湍流封闭模式}\addcontentsline{toc}{section}{9.4 湍流封闭模式}

诺应力(二阶矩)间的关系式之后，雷诺方程和雷诺应力输运方程就封闭了，这种模式 称为二阶矩模式.此外，依据封闭关系式的数学方程形式，湍流模式义可分为代数模 式或微分模式.本 $\tau$\textasciitilde{} 只介绍一些常用的低阶 H 闭模式. 和建立粘性流体的本构方程相仿，构造湍流坷闭模式也需要遵循张量 一 致性、客 观性等基本原则.此外，湍流统计矩还必须满足物理和数学的不等式.如 z

\[<u\tilde{u}~>>O ,  <u\tilde{u}~>>O ,  <U\tilde{iU}\tilde{i}>>O\]

和

\[<u:u:> ~ [<U:U:><U';U:>JI2\]

上式可由数学中施瓦茨不等式导出. 1.布西内斯克( Bus$\tilde{ine}$$\tilde{qe}$) 涡团粘度模式 布阿内斯克(1 87 7>涡闭粘度模式是一种比拟的思想.把湍流微团的脉动比拟为 分子运动的涨落;微团的平均速度比拟为分子的宏观平均速度;分子运动涨落产生的 统计平均动量输运比拟为湍流脉动动世输运.在分子运动的统计理论中分子运动涨 落产生的统计平均动量输运等于宏观的粘性应力;在湍流统计理论中.平均湍流脉动 动最输运等于雷诺应力.因此.根据比拟的思想，由湍流脉动产生的雷诺应力封闭关 系式和分子运动产生的粘性应力有类似的形式.在第 8 章中，已经导出牛顿流体的偏 应力张世有以下的本构方程 z

\[r\upsilon =2\mu S i)\]

根据上述比拟的思想，雷诺应力 \_ <u:u;> 比拟为偏应力町，平均运动 <U ，) 比拟为宏观 速度 U" 于是在以下的不可压缩湍流运动的涡团粘度公式 \_ <u:u:> = 2 川。〉 -3崎 q (9.63) 这里 $\nu$T 称为涡团粘度. < 5'1> = [J<Ui > /cJ.:rj +a<u, > /cJ.:r .]/2 为平均切变率张量，走 为揣动能.上式最后一项是为了满足不可压.缩流体的连续方程，因为式 (9.63) 做张量 收缩时 ， <5 ii > = 0 ，式 (9.63) 最后 一 项使等式得以成立.虽然湍流雷诺应力的涡团粘 度表达式 (9.63) 与牛顿型流体的本构关系式( 9. 的形式上相同，但是 $\nu$T 不是介质的 物性常数，它与湍流的平均场有关. 下面在简单的平衡湍流中 (Pk\textasciitilde{}$\epsilon$) ，通过量纲分析导出川的 一 般表达式.首先问 应是湍流脉动速度量和脉动尺度的乘帜 z

\[\nu r \tilde{qL} (9. 64a)\]

式'11 q = < u:u:> 1/2 = .，f2f是脉动速度的特征址 ， L 为脉动的特征尺度.代入式 (9.63) ， 可得

\begin{equation*}- < u:u '; > - qL < 5i, >\tag{9.64b}\end{equation*}

进一步由局部平衡关系 Pk$\tilde{E}$ 可导出 $\rightarrow$〈。;〉 (EjEi+ 非)三 $\nu$ 〈拦拦〉 按照揣动能的串级理论，揣动能耗散等于大尺度脉动传递的能量.因此，揣动能起散 率 E 可以用揣动能和脉动尺度估计如下 z 民主二=旦: (9.65) 1 [ 将式 (9.64b) 和式 (9. 65) 代入局部平衡关系，就有: ql 侣。 )<Sij)$\tilde{q}$3j[. 简化后得

\[q2 - [2 < Sij ) < S i )\]

则涡团粘度的一般形式 $\nu$， -q[ 可以写作:

\begin{equation*}\nu  T - [2 ( < Sij ) < Si ) ) I 2\tag{9.66}\end{equation*}

式 (9.66) 说明只要得到脉动尺度的表达式，就能得涡团粘度的关系式，这种简单的代 数形式的涡粘模式称为普朗特据合长度模式. [称作混合氏度.湍流的实验数据表 明，脉动尺度与平均切变场有关，不同的平均切变场中. [具有不同的分布规律，例如 在固壁附近 z

\begin{equation*}[ = ky\tag{9.67a}\end{equation*}

y 为距固壁面的垂直距离，走 =0.4-0.41 称为卡门常数，在揣射流或其他无界的自 由湍流切变层中 z

\begin{equation*}[-8\tag{9.67b}\end{equation*}

S 为自由切变层的特征厚度. 2. 瑞动能-耗散率模式(走 -${\rm e}$ 模式) 涡团粘度模式的关键是确定合理的湍流脉动长度尺度[.在简单涡团粘度模式 中 1 由实验结果作最佳拟合得到它的代数式，如式 (9.67a) 和式 (9.67b). 它们显然 是经验性的.很难推广到比较复杂的平均切变流中.事实上脉动长度尺度是各种脉 动成分尺度的平均值，通过对湍流脉动能量输运的机制的研究，可以对 t 做更好的 近似公式.前面已经阐述过，湍流耗散主要在小尺度脉动，或者反过来说小尺度脉 动占有绝大部分的艳散率 z 另一方面揣动能的输入主要来自平均场流动，因此属大 尺度脉动，或者说大尺度脉动占有揣动能的绝大部分.根据揣动能的串级理论和局 部能量平衡近似说，湍流平均脉动长度尺度可以由 k( 揣动能)与 $\epsilon$( 揣能起散率)来 估计.即

\[[ - k:l/2 j\epsilon \]

因此一般涡团粘度为 $\nu$， OCq[ OC k 2 !E.或写成 z

\[\nu  T = CI' k\]

2

\begin{equation*}j \epsilon \tag{9.68}\end{equation*}

其中 C" 为待定常数.用当地揣动能和揣动能起散来表示的涡团粘度模式称为k-E 模 式，为使式 (9.68) 封闭，应当有 h 和 $\epsilon$ 的输运方程.揣动能 h 的方程式 (9. 55) 已在前 面导巾，推导耗散率 $\epsilon$ 的演化方程是较为繁琐的演算，我们给出它的常用形式而略去 推导过程 z ($\tilde{E}$ + (U ，) 主= p, +D, +E, rJ t . . - J ' Jt ${\rm ‘}$e; ${\rm •}$ -.: (9.69) 方程左边是沿平均运动轨迹揣能耗散率的增长率，方程右边的 三 项 P. ，丘，旦分别称 为耗散率的生成、扩散和"耗散"，这些项本身含有高阶的统计量，根据大量的实验数 据和 一 些理论上的考虑可以分别写出这些高阶统计量常用的经验关系式.最后k-E 模式的封闭方程如下 z EE 一+aq ++ AR ----旦、叫主 h 吼叫++ 品

\[-h\]

占

\[-h\]

 (U ，> P. =-(u:u:> 气一

\begin{equation*}r/I,\tag{9.70a}\end{equation*}

D.= 主 I ($\nu$+$\nu$I )乒 1

\[dIkL <I I. ..J\]

(9.70b) E a(u,> P. =- c., (u:u';) 二一一一-${\rm •}$ -., ' -'-J' k JIj (9.70c) Da=52[\{$\nu$+ 叭)去] (9.70d)

\[E. = C.2 (9.70e)\]

将以上方程与涡团粘度方程 z$\nu$T= 【 :"k 2 /$\epsilon$( 式 (9.68)\} 及雷诺平均方程 (9.49a) 、连 续方程 (9.49b) 联立组成封闭方程组.封闭方程 (9.70) 含有常数 C" , Cd . (，: 2 响，它 们需要用典型流动的实验结果和算例结果作最佳拟合来得到.目前常用的经验系 数为 \{飞= 0.09 民= l. 3 (9.70[)

\[C,1 = 1. 45 , C.z = l. 90\]

有了雷诺应力的封闭方程，可以通过求解雷诺平均方程得到平均流场.平均流场的边 界条件提法和层流运动相同，不再重复.对于简单湍流运动，可以利用湍流模式得到 解析解，而复杂湍流运动需要应用计算机求数值解，这是目前工程中常用的预测复杂 湍流运动的方法.

\section*{9.5  圆管中的定常湍流}\addcontentsline{toc}{section}{9.5 圆管中的定常湍流}

下面应用简单涡团粘度模式预测光滑圆管中充分发展湍流的平均特性. 1.圆管瑞流的基本方程 光滑图管中充分发展湍流的平均流动有如下特点. (1)平均速度场是定常平行流 p (2) 除压强外.一切平均提只和困管巾径向坐标有关. 即管内平均流动具有以下性质 z

\begin{gather*}
(U ,), (UO) , (U r ) = O , O, U(r)\\
iJRi, - iJRij - ( 
\end{gather*}

iJ.r iJ()

\[(P) = P(r ,.r)\]

其中 r ，()，.r 为柱坐标，如图 9.8 所示. (9.71a) (9.71 b) (9.71c) (>-4 fL--仨\textasciitilde{} \textbackslash\{\}0) 图9.8 阴管湍流流动 i1主意图 根据以上特性，柱坐标中同管湍流的雷诺平均方程简化为 1 iJ P 1 d , , , '", I d2 U , 1 dU \textbackslash\{\} 一一- -一一 \textasciitilde{}(r(u$\tilde{u}$'，)) + II\{ 一一+ \_\_:\_ \textasciitilde{}.\textasciitilde{} 1 (9.72a) p iJ.r r dr $\nu$ 飞 dr2 ' r dr I 1 iJ P 1 d , , ,.,,, , (u'\textasciitilde{}) $\tau$- -一\_\_:\_ \textasciitilde{}(r(u';)) + 一\textasciitilde{} (9.72b) p r1 r r ar r 阅管壁面上所有速度分量等于零，包括平均速度和脉动速度， 1坷此方程 (9.72 川，方程 (9.72b) 的边界条件为

\[r = R:U = (u:) = (u:u' ,) (u'P = 0 (9. 72c)\]

在代人涡闭粘度公式前，先将式 (9. 72a) 和式 (9.72b) 做适当简化，积分式 (9.72b) 得  rr I ( u';) - (u'\textasciitilde{} ) \textbackslash\{\} P 〈 -ur)+p 〈 u'! 〉十叫 III ,- r' r ,- 11' )dr = Pw( .r) (9.73) 这里 R 为图管半径， P即为回管壁面压强，它只是 r 的函数，由于性质 (2) ，上式中左边 除 P(.r ， r) 外，其余项都与 r 元关，因此有 矶-由叮叮-h (9. 74)

代入式 (9. 72a) 后得 1 dP ". 1 d , , , '" I I d 2 U I 1 dU 、一一一=一一 7(r<U$\tilde{U}$'r>) + $\leftarrow$一+一\textasciitilde{}，\textasciitilde{} 1 (9.75) p d.r ---;: dr\textbackslash\{\} r\textbackslash\{\} uzu r/\} TIJ 飞 dr2 ' r dr J dU 上式对 r 积分一次，由于流动的轴对称性，在轴线上 (U$\tilde{U}$: > (O) =$\gamma$(0) =0 ，于是ar r dP w , , ,,' dU 一-: w =-p(u'rU\textasciitilde{}> + $\mu$$\gamma$(9. 76a) 2 d.r f" -r-z" t" dr 在圆管壁面 (r=R) 上: (u$\tilde{u}$: > = 0 ，而 $\mu$ dU/dr=Tw ' 它等于壁面剪应力，因而进一 步有 dP四 4Tw d.r D 式中 D=2R 为回管直径.将式 (9.76b) 代人式 (9.76 剖，又可得 2芳-一 $\mu$U:> + $\mu$ 忑dU (9.76b) (9. 76c) 现在可以看到，只要有恰当的雷诺应力模式，积分式 (9. 76c\} 就可以求出速度分 布及压力梯度.其计算过程和第 8 章导出哈根-泊肃叶流动的相应公式相仿. 为了反映壁面附近的湍流性质，通常将坐标原点设在壁面，令

\[y = R-r (9.7 7)\]

相应地叫 =-4. 在壁面附近，湍流脉动和平均速度常用壁面摩擦速度无量纲化，它 的定义是 2 1 T". 1 u\textasciitilde{} = 一一一 (9.78) p 在回管流动中压强梯度 dPw/d.r <O ，也就是壁面剪应力 Tw<O ， 因而 u\textasciitilde{} = 1 T".I / p= -Tw/p. 将式 (9.77) 和式 (9.78) 代入方程 (9. 76c) ，得 (U:U'v> 一 $\nu$ 亚=- \_!\_(1- \textasciitilde{}${\rm i}$Tw =- (1 一旦 )U\textasciitilde{} (9.79a) dy p 飞 R \}'W 飞 R J 式 (9.79a) 左边表示流体总剪应力，即雷诺应力和分子粘性应力之和.公式右边表示 总剪应力由壁面线性地减少到轴线上的零值.将以上方程用壁面特征量 Ur ， V 来无量 纲化后，得

\[(U'rU'.> dU.\]

一」7L 一丁一=豆 -1

\begin{equation*}u; ay ,\tag{9.79b}\end{equation*}

其中 U t = 旦 U r (9.80a)

\[J-r\]

y , = 二一- (9.80b)

\begin{equation*}y-pu -y\tag{9.80c}\end{equation*}

y t 称为壁面无量纲坐标，因为它是用壁面特征量 u ， 和 $\nu$ 无 量 纲化的.由式 (9.80b) 可 以导出壁面无量纲坐标 b 和普通无 量 纲坐标 y 之间有以下关系: 主 I 生旦旦>> 1 (9.81)

\[y\nu \]

在高雷诺数回管湍流中 (Re=U"， D / $\nu$ > 10 1 ) .实验数据表明 u ， 坦 (0. 03\textasciitilde{}0. 045)U刑' 因而 如 / y =u,R / 11= u,Re / (2U.. )\textasciitilde{}(O. 015 一 0.0225)Re >> 1.对于给定 y 值，壁面无 量纲坐标 y + 远远大于普通无量纲坐标 y 值.例如: Re= 10 " . 在管中心处 z 豆 =1. 而 Y + \textasciitilde{}2000; 在贴近壁面处 :y=O.O l.而 Y I 句 20. 下面导出圆管湍流的平均速度分布. 2. 圆管瑞流的近壁分层模型 从流向平均动量方程 (9. 79a) 可以看到 z 回管流动中的总平均剪应力(分子粘性 应力 $\mu$ dU / dy 和雷诺应力 -p ( u : u :> 之和)等于壁面摩擦应力和 (1 - y / R) 的乘积.在 靠近壁面的薄层中，豆 = y/R << 1.1- y / R$\tilde{l}$. 这时总剪应力近似等于壁面剪应力，称 这一近壁层为近壁等剪应力层，简称等剪应力层.在等剪应力层中， $\mu$ZE-p 〈心 ' y ) = 凡当 y = y / H << 1 oy 通常等剪应力层近似可以达到 y$\tilde{O}$. 2\textasciitilde{}0. 3. 在高雷诺数的实际流动中该处平均流 速几乎达到管内的最大平均速度.因此等剪应力层内的平均速度可以近似圆管湍流 的平均速度分布. 等剪应力层还可以分成近壁线性次层和对数层. (1)线性次层 在紧贴壁面处(.)1$\rightarrow$ 0) 湍流脉动很小，因此 (u$\tilde{u}$: )句 O. 这时式 (9.79b) 可简化为 利用边界条件 U + (0)=0 ， 得 dU t 唱 dy  $\cdot $ (9.82)

\begin{equation*}U+ = y+\tag{9.83}\end{equation*}

就是说在近壁区湍流的平均速度是线性分布.大量实验结果表明，近壁线性分布的平 均速度在旷 <5 的范围里是很好近似.这一速度分布在形式上与层流状态相似，过 去曾称旷 <5 的近壁区为层流次层，事实上，近代湍流实验证实，这里仍有较强的揣 流脉动，现在称该区为线性底层. (2) 对数层 在线性次层以外，在 ${\rm y}$ <<l. 而 yl >>1 的区域中，分子粘性作用几乎可以忽略，就 是说总剪应力中雷诺应力远远大于分子粘性应力，或者说，湍流涡粘系数远远大于分 子粘性系数，这一区域称为等雷诺应力层.忽略分子粘性应力 dU ， /dy ， 后，

式 (9.79a) 可简化为

\begin{equation*}- (u:u' ,> = u~\tag{9.84}\end{equation*}

应用混合长度模式 (9.66): dU \_ 2..21 dU 1 - (U$\tilde{u}$\textasciitilde{}> = $\nu$ 一一， $\nu$，一，，" y" 1 一::'1T dy' '" - ,. y I dy I 将它代人式 (9.84) 中的雷诺应力项，得平均速度方程如下 z ,,2 y2 ( 艺) 2 = u; (9.85) 用 U. ， II 无量纲化后，得 ? I dU + \textbackslash\{\} 2 ${\rm •}$ ....... dU" kzy1i7-L l =1 , 或 T-L= 一一 飞 ay I I ay+ "y + 积分上式就可以得到圆管中的平均速度的对数分布 z

\begin{equation*}U \phi  =ilny -+ B\tag{9.86}\end{equation*}

K 实测圆管速度证实对数型速度分布是很好的近似，可适用于厕管中 yl >30 , 豆 <0.3 的流动区域，并由实验结果确定式 (9.86) 的常数 z

\begin{equation*}" = 0.4 , B = 5.5\tag{9.87}\end{equation*}

其中 K 称为卡门常数. (3) 缓冲区 5<y + <30 的范围称为缓冲区，这里速度分布既非线性也不是对数型，可以将涡 团粘度和分子粘度公式一起代人式 (9.79) ，用积分求出平均速度分布.实用上，常用 实验给出近似的表达式 z

\[U f = 51ny+-3.05\]

由周管壁面开始的湍流平均速度的线性分布到 U 对数层的分布是壁面附近湍流的特性，在平面 槽道湍流以及压力梯度不大的湍流边界层的近 壁区中也有同样的流动特性，所以上述圆管壁 面附近的湍流分层特性属于近壁湍流的一般特 征，如图 9.9 所示. (4) 中心区 y>(0.3-0. 的称为中心区，这里壁面的 影响十分微弱，因此不能采用壁面特征量.其特 征长度应是半径 R ，元量纲坐标变量应为豆，为 了和对数区的速度分布衔接，常用表达式为 (9.88) 线!缓: U' =2.5Y,+5.5\_ 性 l 冲: \textasciitilde{} 区 1 区 I\textasciitilde{} 10 100 图 9.9 近壁湍流速度分布 \textbackslash\{\}' 争

\[Um.. - U(y) :L:... =-A'lnj+B'\]

u (9.89)

式中经验常数可采用 B* =0. 8,A' =2.44. (5) 截面平均速度 利用分层的速度分布在截面上积分，可以计算截面平均速度 Um ${\rm •}$ 实用上，可以直 接用对数层的速度分布积分计算截面平均速度.因为对数函数在 y=O 处的奇点是 可积的，而缓冲层以下的流量微乎其微;接近中心区对数层的速度分布和中心层的速 度相差也很小.将对数速度分布在圆管截面上积分求和，可得截面平均速度公式: 旦旦= 2. 51n \textasciitilde{}主+ 1. 75 (9.90) Ur J,I 3. 圆管瑞流的阻力系数 可以用壁面平均剪应力计算流动阻力，也可以用沿管流的平均压降来计算流动 阻力，实用上，往往采用后者.定义以下两个无量纲阻力系数. 定义 9.9 元量纲的壁面平均剪应力称为摩擦系数，用 Cf 表示 z

\begin{gather*}
c rE -f -\\
pU~/2
\end{gather*}

定义 9.10 无量纲的平均压强梯度称为沿程阻力系数，用 A 表示:

\[\lambda D dP",\]

=一一一

\[pU~/2 dx\]

由式 (9. 76b) ，圆管湍流沿程压降和壁面平均剪应力之间有以下关系 z

\[dP", 4\tau U'\]

dx D 因此壁面摩擦系数和沿程阻力系数之间有以下的关系 z

\[\lambda  = 4Cf\]

利用式 (9.76b) 和式 (9.78) ，还可以进 一 步导出沿程阻力系数公式: 或

\[\lambda D dP", - 8u~\]

一一一一-

\begin{gather*}
pU~/2 dx U~\\
Ur - ../\AA 
\end{gather*}

Um 2.;2 将平均速度公式 (9.90) 代入式 (9.93a) 后，得沿程阻力系数

\[A-8ui-89\]

一一一 U\textasciitilde{} (2.51n 孚+ 1. 75 f (9.91a) (9.91b) (9.92) (9. 93a) (9.93b) (9. 93c) ufl\_ U r UmR \_ U r n\_\_ n \_ ./f 上式分母中一一-一一一一\_:-$\tilde{Re}$= Re 一一给定 Re 后，求解隐式代数方

\[\nu  Um \nu  2Um\]

4.;2' 程 (9. 93c) 可以得到阻力系数 A 值，一般需要用迭代方法求出该代数方程的解.

4. 圆管平均速度分布的事函数形式和沿程阻力系数公式 用对数函数的平均速度分布是一种近似方法，用它导出的沿程阻力系数公式需 要迭代求解.在工程计算中连用一种更简便的方法计算平均速度和阻力系数，即幕函 数形式的平均速度分布(有人从理论上论证，它是合理的速度分布). (1)事函数型速度分布为 ;=(云 r (9.94a) 其中指数 n 随 Re 的变化见表 9.1 衰 9.1.

\begin{gather*}
Re 4. X 10' 10" 10" >2. OX 106\\
n 1/ 6 1/7 1/9 1/10
\end{gather*}

费引向 o. Hinzc 蕃.捕流. (2) 事函数型阻力系数公式 采用事函数形式的平均速度分布时，圆管的阻力系数也采用事函数型，例如 z 式 (9.94b) 适用于 Re<10 6 A 一 0.3164 一一

\[Rel /-l\]

5. 壁面物理粗糙度、水力粗糙度和相对粗糙度 (9.94b) 前面给出的回管湍流平均速度分布和阻力系数都是流体流经光滑圆的情况，而 通常使用的水泥输水管道表面是不光滑的:钢铁输水管道在使用一段时间后，由于液 体的腐蚀或剥蚀.管道表面也变得不光滑.不光滑表面的平均突起高度定义为壁面粗 糙度或物理粗糙度(见图 9. 10) ，并用 h 表示. h h n 人工粗糙元

\[A /I1i \wedge  II.J, w v V'\]

不规则面 图 9. 10 壁面粗糙度的示意图 壁面粗糙度和壁面摩擦应力有紧密联系，一般来说，表面越粗糙，壁面摩擦应 力也越大.在水力学手册中可以查阅各种材料的管壁粗糙度.物理粗糙度对壁面摩 擦应力的影响和流体的粘性和流动速度有关，流体的粘度大，粗糙度对壁面摩擦应 力的影响就小 z 流动速度低，粗糙度的影响也小.壁面粗糙度主要对近壁面的湍流平

均流动产生影响，用哇面摩擦速度无量纲化的壁面粗糙度称作水力粗糙度，它的定义 如下: h l =但z (9. 95) 已有的研究结果表明:当 h ' <5 时，壁面粗糙度对摩擦应力没有影响，这种状态称作 水力光滑管.从前面陈述的分层速度分布规律可知，当 11<5 $\nu$ / U r 时，粗糙突起完全淹 没在线性底层内，它对等雷诺应力层的影响很小.因此既不影响对数分布律，也不影 响壁面摩擦应力.因此 $\nu$ <5 的情况可以视作光滑哇面.当 h 1>100 时，粗糙度对流 动的影响达到饱和状态，壁面摩擦应力完全由粗糙度决定.这时，粗糙突起已进入对 数区，它对等雷诺应力层有很大的干扰，光滑阴管的理论完全不适用.粗糙突起是产 生噎面摩擦的主要因素. 水力粗糙度(式 9.95) 中含有壁面摩擦速度，而壁面摩擦速度就是壁面摩擦应力 的另一种表达方式，它是需要计算的量.所以，用水力粗糙度作为参数来建立公式需 要用迭代方法计算阻力系数，实际使用时不方便.丁-程中常用另一种元量纲的粗糙度 表达式，称作相对粗糙度.它的定义是物理粗糙度除以圆管直径，即 I${\rm I}$= 去 (9.96) 考虑粗糙度的阿管湍流是一个比较复杂的湍流问题.从水力摩阻手册可以获得 常用的工程计算公式;也常用著名的尼古拉茨曲线来计算厕管流动的阻力系数. 6. 尼古拉茨曲线 德国工程师尼古拉茨在 1933 年对图管流动的压降作了系统的测量，整理成工程 计算的曲线图〈图 9. 11) 称为尼古拉茨曲线，图中横坐标是流动的雷诺数 Re ， 纵坐标 是阻力系数 $\lambda$ ，相对粗糙度作为参数表示在不同的曲线上. 图 9.11 表示阻力系数随雷诺数和粗糙度的变化.从尼古拉茨曲线可以获得回管 流动的全部阻力特性.

\[(1) Re<2300\]

流动的阻力系数随雷诺数的增加而减小，有倒数关系 z

\[A = 64 --\]

Re 该式在第 8 章已经导出，是因管层流流动的精确解.在层流状态下，壁面粗糙度对层 流状态的阻力系数几乎没有影响.

\[(2) 2300<Re<5000\]

阻力系数突然增加，这是因为流动由层流转变为湍流.在这一雷诺数范围内，阻 力系数没有确定规律，表面剪应力随环境因素(如 z 壁面粗糙度、流动的稳定性等)变 化很大.

\[\lambda 0.1\]

0.020 0.012 0.010

\[(3) Re>5000\]

RI! 图9.11 阴管流动的阻力系数(尼古拉获闹线〉 0.02 0.015 0.01

\[40\omega 8 0.006\]

0.004 0.002 0.001 g:3883 0.0004 0.0002 0.0001

\[0.00\infty 5\]

流动达到完全湍流状态(也称充分发展湍流状态) .阻力系数继续随雷诺数的增 大而减小，不过减小得比较缓慢.充分发展的圆管湍流中，壁面粗糙度对阻力系数的 影响较大.给定雷诺数，相对粗糙度增加，阻力系数增大 z 当相对粗糙度增加到 一 定量 以后，流动进入完全粗糙状态(见图上的虚线) .在完全粗糙区，阻力系数不随雷诺数 变化，而是常数.常数阻力系数表示，压降和管流速度的平方成正比(见阻力系数公 式 (9. 92a)). 图 9.11 中的长虚线给出完全粗糙和 一 般水力粗糙的分界. 7. 非圆截面长直管道的流动阻力 非圆截面长直管道的流动状态和长直圆管流动不同，最显著的流动特性是存在 管截面上的二次流.二次流是一种复杂流动，对于它们的准确预测方法正在研究中. 这里，要强调指出:非圆截面长直管道中的二次流是湍流的特性，非圆截面长直管道 的层流中没有二次流.根据一些实验和理论分析的结果，矩形和 三 角形截面内的湍流 平均流动(二次流)示意于图 9. 12. 非因截面管内湍流的阻力系数的工程估算采用当量水力半径的方法，就是说，把 非图截面用某种方法折算成圆截面来计算阻力系数.当量水力半径 rh 的定义是

\[h~\tilde{I} <?/CJ61\]

段写l:J I \textasciitilde{}\textasciitilde{}吏| (a) (b) 图 9.12 非阅截面管内湍流的 二 次流 (a) 矩形截面中的 二 次流. (b) 气 角形截面内的 二 次流 r. = \textasciitilde{}$\times $过水断面 h 湿润周长 第 9 章 ;湍流 (9.97) 过水断面是管内水流的实际流通面积\$湿润周长是流体和管壁接触的边界长度(见 图 9.13 )，对于充满管道的流动，过水断面等于管道截面权.湿润周长等于截面积的 周长.对于不充满的管道流动，根据实际流通面积和湿润周民计算水力半径.读者可 以验算，对于充满管道的圆管流动，水力半径恰好等于回截面的半径，这是构造水力 半径公式的依据.对于方形截面，或者不是十分扁长的矩形或椭圆形截面，应用水力 半径计算管流阻力是很满意的.至于，扁长截面的管流，水力半径的方法不可取，必要 时，需要应用流体力学的一般原理加以计算. 国|---;;:- 'I-\textasciitilde{}- l$\tilde{OO}$ (a) (b) (c) 图9.13 说明计算水力半径用图 2,," (a) a X a 矩形 .r'=7' (b) a X b 矩形 .r，. = 一一-:; (c) 充液丽权 aXb.r.= 一一一a 十 bii \textbackslash\{\}\textbackslash\{\}.:1 7 L.tr:l.UlIl" " " U"" - 2a + h 本节以上讲述的长直管道流动沿程阻力只是有管壁摩擦产生的阻力，不包括输 液管道中常有的阅门、弯头、突然缩放等部件产生的附加阻力，这些阻力称为局部阻 力.关于局部阻力的讨论不属于本书范围，读者可以查阅工程设计手册得到相应的局 部阻力系数公式和数据.

\section*{9.6  切变湍流的拟序结构}\addcontentsline{toc}{section}{9.6 切变湍流的拟序结构}

本章前几节以统计方法为基础研究湍流的平均运动，这 一 方法导出了不封闭的 方程组，为了封闭它必须对脉动场的高阶统计矩作封闭性假设.合理的封闭模式无疑 地有顿于对脉动场的了解，经典的湍流统计理论曾经对随机分布脉动速度作球涡或

椭球祸的假定，并得到了比较满意的均匀湍流衰减的预测.然而在工程实际中很普遍 的非均匀切变湍流，它们的脉动场要复杂得多.为了对它有深入的了解，近 50 多年来 人们用流动显示和测量方法对切变湍流脉动场进行了大量研究，取得了丰富的实验 数据，发现了切变湍流中存在拟序流动结构. 所谓拟序结构(或称相干结构)是指湍流脉动场中存在某种序列的大尺度运动， 它们在湍流场中触发的时间和地点是不确定的，但一经触发就以某种确定的序列发 展.这一发现是湍流研究中重大的突破，因为它表明湍流并不像经典理论想象的那样 是完全不规则的，而是在小尺度不规则的背景脉动中存在若干有序大尺度运动.下面 介绍两种典型的拟序结构. 1.壁瑞流的拟序结构 在 9. 5 节中得到了壁面附近的平均速度场，那里存在线性次层、缓冲层和对数 层.自 20 世纪 50 年代末以来，不少湍流研究人员对近壁湍流脉动进行深入研究，由 于这里流速较低，比较容易应用染色线方法来直接观察流动状态.常用的显示方法有 注入有色液体线，然后追踪染色线的变形来观测流场结构.比较成功的方法是氢泡线 法.将脉冲电压施加于极细 (5\textasciitilde{}10$\mu$m) 的金属阴极丝上，产生脉冲电解氢气泡线.跟 踪氢泡染色线可以推断流场的脉动形态.后面提供的拟序结构的显示图像大部分是 用氢泡线方法显示的.实验结果发现在线性层和缓冲层中存在如下的拟序结构. 图9.14 近 M 条件?结构的氢泡显 ${\rm i}$f (a 1 条带纺构示意网 ， (b> 条带结构照片

的条带结构.条带的平均间距 l = 100 $\nu$ / 饵，条带长短不 一 .平均长度为 l = 1000$\nu$ / 帆$\cdot $ 从氢气泡线上可以判断，条带处于低速区(条带处的速度小于当地平均速度) ，进 一 步 分析可以断定低速条带是流向涡的痕迹. (2) 条带的升起、振动和破裂 条带 一 经出现，便开始缓慢升起(见图 9. 15) ，约在 y l= 15-30 处发生振动后突 然破裂.从涡的运动学性质及显示图像来看，上升的条带是 一 种 H 形或马蹄形涡 (图 9. 15). 在上升的马蹄涡头部区域流向速度剖面上出现拐点，并形成局部高切变 率层，这种速度剖面是极不稳定的.于是升高的条带(或马蹄涡)发生振动，随后很快 破裂.条带突然破裂的过程又称猝发.条带猝发过程产生极大的动 量 输运'瞬时的动 量 通量最大值 ($\omega$u 、$\gamma$')n阳川a 个猝发过程的平均脉动动量通量约为系综平均值的 7刊0\% 左右.就是说，在湍流边界 层内层，雷诺应力是在短时间内由猝发产生. -- U<z .t )l 气泡线) 图9.15 上升条带的振动及破裂 (3) u 下扫"和条带的再现 条带的破裂伴随有一股强烈的流向加速和向下的流动，加速向下流动称为"下 扫飞下扫过程的平均脉动动 量 通 量 约为系综平均值的 30\% 左右.下扫的扰动在壁面 附近能再次诱发新的条带结构，于是又 一 次出现拟序运动.并不是所有"下扫"运动都 诱发条带，另外条带的持续时间也有长有短，都受环境的影响.现有的实验结果表明 在底层出现拟序结构的平均周期或两次猝发的平均时间间隔约为 TJj = 6 8/ U时 ， ${\rm o}$ 为 边界层厚度. 湍流边界层近壁拟序结构的发现充分说明，湍流拟序结构是维持湍流，即产生雷 诺应力的主要机制.我们可以设想，如果能控制近壁的拟序结构就有可能控制近壁揣 流，从而达到减阻的效果.事实上，已经有 一 些实验装置实现近壁结构的控制，例如， 带细掏槽的壁面、壁面吸气、壁面温度控制以及声控等. 上面描述的近壁拟序结构是极初步的，由于拟序结构对湍流生成有重要作用，对 它们的细节和动力学机制正在进行更深入的研究.

\section*{习题}\addcontentsline{toc}{section}{习题}

2. 混合层的拟序结构 不受固壁影响的自由切变层中也存在拟序结构，布朗和劳照柯 (Brown- Roshko) 首次得到混合层中拟序结构的流动显示图像.混合层是由两股速度方向相同，大小不 等的流动汇合生成.为了防止由密度分层引起的其他不稳定性，混合层的一侧是氨和 氧的混合气体，另一侧是密度相同的氨气.由显示图中可以看到棍合层开始有一段层 流不稳定向湍流过渡的过程，而后在充分发展的湍流区仍有可辨认的大涡结构，在大 涡结构上附着小尺度湍流.图 9.16 是典型的混合层中拟序结构的图像. 图9.16 混合层的拟序结构(流动由左向右) 人们发现自由切变层中的大涡结构是产生噪声的主要来酶，遏制大涡能够减少 噪声.

\noindent\textbf{9.1.} 导出柱坐标系中不可压流体平面湍流流动的平均流连续方程和雷诺方程-

\noindent\textbf{9.2.} 证明不可压缩牛顿流体湍流平均运动的动能满足如下方程 (U; 是平均速

度 ， U ; 为脉动速度，质量力不计) : a ,U;U; \textbackslash\{\} 'TT a lUiUi 飞子 r\textasciitilde{}'\textasciitilde{}\textasciitilde{}'I+U; \_u一 1 \textasciitilde{}'\textasciitilde{}\textasciitilde{}' I dt 飞;:: J . rJ Xj 飞$\tilde{I}$ r U;P , iJ ,U,U i \textbackslash\{\} TT' I ,, 1 iJUj iJUi - (u\textasciitilde{} 叫〉二二+一一 I- ::::L:.... +$\nu$ l 一$\tau$一卜 U，〈 u d UJ| - v 一一一- l' dX${\rm i}$ ' cJ X J L p 0 - dXj 飞 2 I -, o\_-,--, oJ -axJ${\rm Ð}$xJ

\noindent\textbf{9.3.} 证明二维湍流巾，脉动流动的能量耗散函数是

.p =$\mu$(21 去 1 2 +2 时 1 2 + 1 努 1 2 + 1 主 r +21 苦苦 1 )

\noindent\textbf{9.4.} 直径为 50mm 的铜管道内水流平均速度为1. 7m/s. 测得壁面剪应力为

\[6. 724N/m\]

2 ${\rm •}$ 已知 $\nu$= 1. 006mm 2 / s. 试计算粘性次层、过渡层及完全湍流层的厚度. 9.S 在 Re<10 的条件下，若已知光滑阅管内湍流的平均流速度分布为

;=(l 一云) 1/7 式中 r 为离轴线的距离 .R 为圆管半径，请导出下列等式 z 49 (1)管的平均速度为 Um=;$\tilde{U}$60 (2) 混合长度 L=7R 卫!\_ ( 1 \_ \textasciitilde{} ) 617 ( \textasciitilde{} ) 1/ 2 Uma 嚣 \textbackslash\{\} $\tilde{R}$ J 飞 R J

\noindent\textbf{9.6.} 已知光滑管速度分布为

U n" ${\rm •}$ r (R - r) Ur ll / 7 二= 8.741 一一一一一一| U. L I! \_j 0.305 \_..L...., 2Um R 请导出 z$\lambda$=$\tau$节.，式中 Re= 一一一.

\[\tilde{e} -' - I!\]

\section*{思考题}\addcontentsline{toc}{section}{思考题}

第 9 章湍流 S9.1 通过坐标变换，雷诺应力张量必能变换为三个大小不等的正应力，其他分 量为零，为什么? S9.2 为什么说任何固定壁面上的揣动能和雷诺应力都等于零?揣动能生成率 和超散率也等于零吗?为什么? S9.3 在槽道或回管中的湍流运动，中心平均速度最大，该处的"湍流脉动和揣 功能生成率也最大"的说法是否正确?为什么? S9.4 应用混合长度模式，能否估计湍流粘度和分子粘性的量级之比?
